Qubits are scarce, so quantum software has to treat temporary workspace as a
real resource. A qubit that has finished one job may be cleaned and reused for
the next. The practical question sounds simple: when is that reuse actually
safe? This paper argues that the answer has two parts. The qubit must be in a
known, separable state, and the system must also prove that no hidden
information flow survives through it.

The first case study gives the constructive side. In coherent teleportation,
deferred measurement replaces classical feed-forward with controlled quantum
gates. The circuit factors as
$\ket{++}\otimes\ket{\psi}$, and two final Hadamards produce the exact
postcondition $\ket{00}\otimes\ket{\psi}$. The source and ancilla are therefore
known, separable, and reusable in the ideal model.

The second case study gives the adversarial side. Exact partial traces of three
educational spy-detector exports show that the advanced circuit reproduces the
baseline computational-basis distribution on the displayed four-wire register
to within $8.90\times10^{-9}$ total-variation distance, while its purity falls
from $1$ to $1/2$, its fidelity with the baseline is $0.4268$, and its trace
distance is $0.6184$. A new fixed-input, branch-conditioned reconstruction then
resolves the retained-wire semantics. The routing network maps
$(q_2,q_3,q_4)$ to $(q_3,q_4,q_2)$, placing the spy-side measurement branch on
$q_4$. For Alice's value $v$, basis $b$, and the spy's basis $e$, the averaged
fifth-wire state is $\ket{v}\!\bra{v}$ when $e=b$ and $I/2$ when $e\neq b$.
Within this reconstructed model, the legitimate subsystem is preserved with
unit fidelity to floating-point precision.

The result motivates a broader threat model: an untrusted compiler, runtime,
source, or control stack may use reclaimed qubits as hidden workspace while
preserving only the outputs a user verifies. We prove the corresponding
information-disturbance boundary, distinguish four levels of circuit
equivalence, and introduce a hyperlinked 5-by-6 security capability matrix that
separates feasible, conditional, hypothetical, blocked, and unnecessary attack
objectives across five adversary-access tiers. A linked appendix records what is
touched, what may remain untouched, available alternatives, mitigations, and
evidence for every cell. We also release a versioned research artifact with
automated tests and deterministic regression checks. The exact Boolean acceptance register of the
original educational detector, full Qiskit parity, and hardware validation
remain explicit open problems.
